\titledquestion{Une autre d\'efinition du mouvement brownien}
\begin{parts}
    % @QUESTION
    \part Montrer qu'un processus $(X_t)_{t\in \R^+}$ est un mouvement brownien si et seulement si:
        \begin{itemize}
        \item[i')] $X_0=0$,
        \item[ii')] $(X_t)_{t\in \R^+}$ est un processus gaussien centr\'e,
        \item[iii')] ${\rm Cov}(X_s,X_t)=\inf(s,t).$
        \end{itemize}
    % @QUESTION_END
    \begin{xsolution}
    % @SET_START
    Commençons par le sens : brownien $\Rightarrow$ processus gaussien.
    % @LIST_START
    \begin{enumerate}
        \item Structure d'indépendance des accroissements :  
           % @ATOM
           % @PRECOND
           Pour \(s \leq t\), on décompose \(X_t = X_s + (X_t - X_s)\).  
           % @ARGUMENT
           L'indépendance de \(X_s\) et \(X_t - X_s\) 
           % @ARGUMENT_END
           implique que la covariance se réduit à :  
           \[
           \mathrm{Cov}(X_s, X_t) = \mathrm{Cov}(X_s, X_s) + \mathrm{Cov}(X_s, X_t - X_s) = \mathrm{Var}(X_s) + 0 = s
           \]  
           ce qui donne directement 
           % @OUTCOME
           \(\mathrm{Cov}(X_s, X_t) = \min(s,t)\).
           % @ATOM_END
        \item Caractère gaussien :  
           % @ATOM
           % @PRECOND
           Le brownien standard a des accroissements gaussiens et 
           % @PRECOND
           \(X_0 = 0\), 
           % @PRECOND_END
           donc toute combinaison linéaire finie de \((X_{t_1}, \dots, X_{t_n})\) est 
           % @ARGUMENT
           une combinaison linéaire d'accroissements gaussiens indépendants, elle est donc gaussienne.
           % @OUTCOME
           Le processus est donc gaussien centré.
           % @ATOM_END
    \end{enumerate}    
    %Ces deux points exploitent essentiellement \emph{l'indépendance des accroissements} et la \emph{normalité} des accroissements.
    % @LIST_END
    Puis la réciproque:
    % @LIST_START
    \begin{enumerate}
    \item Calcul de covariance des accroissements : 
       % @ATOM
       % @PRECOND
       En utilisant la propriété ${\rm Cov}(X_s,X_t)=\inf(s,t).$
       % @PRECOND
       Calculer $\text{Cov}(X_t - X_s, X_v - X_u)$
       Pour \(s < t \leq u < v\) :
       % @ARGUMENT:CALCUL
       \[
       \text{Cov}(X_t - X_s, X_v - X_u) = \min(t,v) - \min(t,u) - \min(s,v) + \min(s,u)
       \]
       Ceci donne systématiquement 0, prouvant que 
       % @OUTCOME
       les accroissements disjoints sont non corrélés, donc indépendants 
       % @OUTCOME_END
       (car vecteur gaussien).
       % @ATOM_END
    \item Calcul de la variance des accroissements :  
       % @ATOM
       % @PRECOND
       Calculez la variance d'un accroissement de $X$:
       % @PRECOND_END
       Pour \(s < t\) :
       \[
       \text{Var}(X_t - X_s) = t + s - 2s = t - s
       \]
       % @ARGUMENT
       donc la loi ne dépend que de \(t-s\).
       % @OUTCOME
       Ceci montre que les accroissements sont stationnaires et 
       % @OUTCOME
       de loi \(\mathcal{N}(0, t-s)\).
       % @ATOM_END
    
    Ces deux points 
    % @SET_START
    % @ATOM
    % @PRECOND
    les incréments disjoints sont indépendants, et 
    % @PRECOND
    la loi des accroissements est gaussienne centrée de variance $t-s$, c'est-à-dire \(\mathcal{N}(0, t-s)\),
    % @ARGUMENT
    plus l'hypothèse de continuité des trajectoires 
    % @ARGUMENT_END
    donnent 
    % @OUTCOME
    la définition classique du mouvement brownien.
    % @ATOM
    % @PRECOND
    les incréments disjoints sont indépendants, et 
    % @PRECOND
    la loi des accroissements est gaussienne centrée de variance $t-s$, c'est-à-dire \(\mathcal{N}(0, t-s)\),
    % @ARGUMENT
    en prenant la version continue via le théorème de Kolmogorov
    % @ARGUMENT_END
    donnent 
    % @OUTCOME
    la définition classique du mouvement brownien.
    % @SET_END
    \end{enumerate}
    % @LIST_END
    % @SET_END
    \end{xsolution}

    % @QUESTION
    \part Montrer que si $W_t$ est un mouvement brownien, $ \{X_t=t\;W_{\frac
    1t}, t>0, X_0=0 \}$
    l'est aussi.
    \begin{xsolution}
    % @SET_START
    % @ATOM
    La Continuité en 0,
    % @PRECOND
    calculer $\lim_{t \to 0^+} t W_{1/t}$:
    % @PRECOND_END
   \[
   \lim_{t \to 0^+} t W_{1/t} = \lim_{s \to \infty} \frac{W_s}{s} = 0 \quad \text{p.s.}
   \]
   % @ARGUMENT
   par la loi des grands nombres pour le brownien.
   % @OUTCOME 
   $\lim_{t \to 0^+} t W_{1/t} = 0$ a.s.
   % @ATOM
   \\
    Le caractère gaussien centré :  
    % @PRECOND
    Pour tout \(t_1, \dots, t_n\), le vecteur \((X_{t_1}, \dots, X_{t_n})\) 
    % @ARGUMENT
    est une combinaison linéaire des \(W_{1/t_i}\)
    % @OUTCOME
    donc il est gaussien centré.
    \\
    % @ATOM
    La covariance brownienne : 
    % @ARGUMENT:CALCUL 
       \[
       \mathbb{E}[X_t X_s] = ts \cdot \min\left(\frac{1}{t}, \frac{1}{s}\right) = \min(t,s)
       \]
    % @OUTCOME 
    i.e. $\mathbb{E}[X_t X_s] = \min(t,s)$.
    \\
    % @ATOM
    % @PRECOND
    continuité: $\lim_{t \to 0^+} X_t = 0$ a.s.,
    % @PRECOND
    gaussianité: le vecteur \((X_{t_1}, \dots, X_{t_n})\) est gaussien centré pour tout ensemble de temps, et
    % @PRECOND
    covariance: $\mathbb{E}[X_t X_s] = \min(t,s)$,
    % @PRECOND
    avec \(X_0 = 0\)
    % @ARGUMENT
    ce qui caractérisent le mouvement brownien.
    % @ARGUMENT_END
    Ces trois points (continuité, gaussianité, covariance) suffisent à conclure que 
    % @OUTCOME
    \(X\) est un mouvement brownien.
    % @SET_END
    \end{xsolution}
\end{parts}